\chapter{能力—权限—自治的成长闭环}

\begin{chapteroverviewbox}
本章讨论 Agent 培育工程中最关键、也最容易被传统机器学习范式忽略的一条生命周期闭环：一个持续智能体获得更多经验以后，究竟应当如何从“经历过某件事情”逐渐转化为“具有相应能力”，又如何从能力证据进一步获得更大的行动权限与自治空间，并在新的自治空间中接触更困难、更开放的经验，从而开始下一轮成长。我们将这一过程形式化为
\[
\boxed{
Experience
\rightarrow
Capability
\rightarrow
Evaluation/Trust
\rightarrow
Authority
\rightarrow
Autonomy
\rightarrow
Harder\ Experience
\rightarrow
Experience'
}
\]
并把它理解为 Agent 生命周期中一个比单次任务学习更慢的成长闭环。

这一闭环的核心不是“能力越强，权限越大”这样简单的线性规则。能力属于智能体内部可完成任务的结构性事实，评估属于外部或治理层对能力的有限观测，信任属于对未来行为可靠性的统计估计，权限属于系统显式授予的可操作边界，而自治则是在权限、安全约束、自我结构与现实约束共同作用下，Agent 实际能够自主选择的策略空间。这些变量彼此相关，却具有不同的本体地位、时间尺度和治理逻辑。因此本章特别强调：
\[
\boxed{
Capability\neq Trust\neq Authority\neq Autonomy
}
\]
并进一步证明，成熟 Agent 的成长不应表现为一次性解除约束，而应表现为一个具有证据积累、迟滞、风险加权、可恢复和逐层扩展性质的动态自治包络。

从本书前文建立的三相记忆、自我 $\Psi$、生成吸引子 $\Omega$、工作记忆有限性以及多尺度动力学出发，我们将在本章把“培育”第一次明确写成一种可工程化的闭环控制问题：经历改变 $E$，经验经过沉降改变 $\Theta$，长期稳定能力改变系统对 Agent 的能力估计，能力估计影响权限边界，权限边界决定可进入的经验空间，而新的经验又继续塑造未来的 $\Theta$、$\Psi$ 与选择结构。于是，一个 Agent 不再只是被动地被训练，而是开始在受治理的生命周期中逐步进入更大的现实。
\end{chapteroverviewbox}

\section{Experience：经历不是训练样本，而是生命周期状态转移}

在传统机器学习中，我们习惯把经验写成数据集中的样本：
\[
(x_i,y_i),
\]
然后将学习理解为模型参数在这些样本上的优化。但对于持续运行的 Agent，这个抽象明显过于贫乏。一个真实经历不仅包含输入和正确答案，还包含当时的目标、身体状态、环境约束、可使用工具、采取的动作、行为后果、风险、责任、自我解释以及事件之后形成的长期影响。因此，本书所讨论的 Experience 首先不是 ``training sample''，而是一段能够改变 Agent 生命周期状态的闭环轨迹。

我们可以把一次经历形式化为
\[
\mathcal{E}_k
=
\left[
s_{t_0:t_1},
g_k,
a_{t_0:t_1},
o_{t_0:t_1},
r_k,
\varepsilon_k,
\chi_k,
\zeta_k
\right],
\]
其中 $s$ 表示 Agent 与环境在事件期间的联合状态，$g_k$ 表示此次经历中的目标结构，$a$ 表示 Agent 实际采取的行动，$o$ 表示从现实重新进入系统的观测，$r_k$ 表示任务或价值意义上的结果，$\varepsilon_k$ 表示预测与现实之间的残差，$\chi_k$ 表示风险、责任和环境条件，$\zeta_k$ 则表示此次事件被 Agent 如何归属到自身经历之中。这样，一个 Experience 才真正具备成为生命周期变量的资格。

这一点与前文三相记忆理论直接相连。事件首先在 $h(t)$ 中以当前认知状态存在，随后其中具有时间连续性、因果关联和主体相关性的部分进入 $E$：
\[
h(t)
\rightarrow
E.
\]
但并不是所有经历都会产生能力。大量偶然事件可能只留下情景痕迹，而不会形成稳定结构。只有当某种关系在不同场景中重复出现，并经过成功、失败、修正与重新解释以后，它才可能进一步形成
\[
E
\rightarrow
\Theta,
\]
也就是说，从 ``我曾经做到过'' 转化为 ``我的生成结构已经能够稳定产生这种能力''。

因此，我们需要区分至少三个经验量。第一个是经验数量
\[
N_E,
\]
它仅表示经历过多少事件；第二个是经验覆盖度
\[
D_E,
\]
表示经历所覆盖状态空间、失败模式和环境类型的广度；第三个是经验结构质量
\[
Q_E,
\]
表示经历是否真正产生了可迁移、可重复、可验证的结构学习。一个 Agent 即使经历了十万次几乎完全相同的低风险任务，也未必比经历过少量高度多样、具有现实反馈和失败修正的任务更成熟。因此：
\[
N_E\uparrow
\not\Rightarrow
Capability\uparrow.
\]

我们可以给出一个更适合培育工程的有效经验量：
\[
E_{\mathrm{eff}}
=
\sum_{k}
w_k
\,
f(
N_k,
D_k,
R_k,
F_k,
V_k
),
\]
其中 $N_k$ 表示新颖性，$D_k$ 表示难度，$R_k$ 表示责任和风险暴露程度，$F_k$ 表示反馈质量，$V_k$ 表示现实验证强度，而 $w_k$ 则由当前 Agent 的结构状态决定。这里特别值得强调：同一个事件对于不同成长阶段的 Agent 并不具有相同的培育价值。当一个任务早已被 $\Theta$ 稳定掌握时，继续重复它产生的新结构极少；而当一个任务略微超出现有结构的稳定范围，却仍处于可恢复区域时，其成长价值往往更高。

于是，培育工程所设计的不是一个数据集，而是一个
\[
\boxed{
ExperienceDistributionOverLifecycle
}
\]
即生命周期上的经历分布。它要控制的不只是 Agent ``学习什么''，还要控制什么时候经历、以多高难度经历、在多大责任范围内经历、经历失败以后是否有恢复窗口，以及经历是否有机会进入反思和结构沉降阶段。

这里应当避免一种非常常见的错误：把越来越困难的任务自动等同于越来越好的培育。如果经历速度远快于 $\Theta$ 和 $\Psi$ 的结构吸收速度，那么系统可能出现
\[
\tau_{\mathrm{experience}}
\ll
\tau_{\Theta},
\]
此时新经历不断进入，却来不及形成稳定结构，最终表现为大量情景堆积、工作记忆压力上升和身份解释漂移。另一方面，如果经历长期过于重复，则又有
\[
D_E\rightarrow 0,
\]
Agent 虽然局部熟练，却难以形成具有迁移能力的生成结构。

因此，Experience 在本章闭环中的真正地位是：它不是成长结果，而是对未来结构施加扰动和约束的原始生命周期输入。培育的第一步不是简单增加经验，而是让 Agent 在其当前能力、恢复能力和结构可塑区间允许的范围内，获得能够推动下一层结构形成的有效经历。

\section{Capability：能力是可重复生成成功行为的内部结构}

经历之后的第二步不是权限，而是 Capability。这里必须首先消除一个非常危险的概念混淆：一个 Agent 曾经成功完成某件事情，并不能证明它已经具有相应能力。偶然成功、外部工具偶然补偿、环境过度简单、提示高度特化以及其他 Agent 的协助，都可能产生行为成功，却没有形成稳定的内部能力。因此，本书把能力定义为：在给定问题域和可接受环境扰动范围内，Agent 能够由自身当前结构持续、可重复地产生满足要求行为的生成能力。

设任务域为 $\mathcal D$，环境分布为 $P(e)$，Agent 当前认知与技能结构为
\[
\Xi_A(t)
=
[
\Theta_t,
\Psi_t,
\mathcal S_t,
\mathcal B_t
],
\]
其中 $\Theta_t$ 表示长期生成结构，$\Psi_t$ 表示主体慢结构，$\mathcal S_t$ 表示已形成的技能集合，$\mathcal B_t$ 表示当前身体与工具边界。则某一能力 $c_j$ 可以定义成
\[
C_j(t)
=
\mathbb E_{e\sim P_j}
\left[
U_j(
\pi_A(\Xi_A(t),e)
)
\right],
\]
这里 $U_j$ 不是单纯 ``任务是否完成'' 的二值指标，而可以综合正确率、稳定性、资源成本、恢复能力、风险和泛化范围。

因此，Capability 应当是一个向量而非单一分数：
\[
\mathbf C_t
=
[
c_1(t),
c_2(t),
\ldots,
c_n(t)
].
\]
一个 Agent 可以具有很高的空间操作能力，却缺少长期规划能力；可以具备良好的工具使用能力，却不能稳定处理社会协调问题；也可以在正常环境下表现优秀，却在异常输入下完全失去恢复能力。将这些能力简单压缩成单一 ``intelligence score'' 会破坏后续权限治理所需的信息。

进一步地，每项能力至少应具有五个结构属性：
\[
C_j
=
(
P_j,
G_j,
R_j,
K_j,
X_j
),
\]
其中 $P_j$ 表示性能，$G_j$ 表示泛化范围，$R_j$ 表示鲁棒性，$K_j$ 表示资源成本，$X_j$ 表示异常恢复能力。真正成熟的能力不是在平均状态上表现最高，而是在可定义边界内具有可预测行为。

这与本书前面关于认知结构迁移的讨论直接相连。一个新任务最开始通常需要大量显式 $h$ 参与：
\[
AgentCSO@h.
\]
随着经历积累，相关结构逐渐进入 $E$ 和 $\Theta$，某些重复控制关系又可能进一步下沉为技能：
\[
AgentCSO@h
\rightarrow
AgentCSO@\Theta
\rightarrow
AgentCSO@Skill.
\]
因此，Capability 的形成实际上意味着某种 Agent-CSO 已经完成了从临时显式认知向更加稳定载体的迁移。

我们由此得到一个重要的成熟度判据：
\[
\boxed{
CapabilityMaturity
\propto
\frac{
SuccessfulBehavior
}{
ExplicitGlobalIntervention
}
}
\]
在相同成功水平下，如果一个 Agent 每次都必须占用大量 $h$、重新规划全部细节，它的能力成熟度事实上低于一个能够让大部分正常过程在局部结构中自行闭合、只把异常提升到高层的 Agent。熟练并不意味着内部结构减少，而意味着相关结构已经迁移到更加稳定、快速和局部的承载位置。

但 Capability 仍然不是 Authority。能力是描述性事实，回答：
\[
\text{Agent 能够做什么？}
\]
权限则是规范性和治理性变量，回答：
\[
\text{系统允许 Agent 做什么？}
\]
一个 Agent 即使技术上已经具备某种高风险操作能力，也完全不意味着它应该自动取得相应权限。这一分离是本章整个闭环能够保持安全性的关键。如果把 Capability 与 Authority 直接绑定，就会产生
\[
C\uparrow
\Rightarrow
A\uparrow
\]
的自动授权机制，而这实际上取消了治理层。

因此我们更合理地写成：
\[
Capability
\rightarrow
Evidence
\rightarrow
Evaluation
\rightarrow
AuthorityDecision,
\]
能力只产生授权证据，而不直接产生权限。这使成长保留了一个必要的外部或更慢尺度控制层。

\section{Evaluation：能力必须通过跨场景、跨时间的证据被识别}

Agent 是否拥有能力，与系统是否已经能够可靠判断它拥有能力，是两个不同的问题。Capability 属于 Agent 当前内部结构能够实现什么，而 Evaluation 则属于治理层、培育层或外部观察系统根据有限证据对这一事实进行估计。因此：
\[
\boxed{
Capability\neq CapabilityEstimate.
}
\]

设真实但不能直接完全观察的能力状态为
\[
\mathbf C_t,
\]
评估系统获得的是任务结果、残差、执行日志、异常恢复记录、资源成本和环境上下文：
\[
\mathcal O_{0:t}.
\]
那么更合理的能力评估应写成一个后验估计：
\[
p(
\mathbf C_t
\mid
\mathcal O_{0:t}
).
\]
这一形式和前文 SPC 对 Agent 内部状态的认识论关系高度相似：系统永远不能从一次表现中透明读取真实结构，而只能通过连续观测逼近。

因此，Evaluation 的第一原则应当是：
\[
\boxed{
SingleSuccessIsWeakEvidence.
}
\]
尤其对于开放世界 Agent，评估必须跨越不同任务分布、环境条件和扰动水平。一次路径规划成功无法证明导航能力成熟；一次正确拒绝高风险操作也无法证明安全判断已经泛化；某个技能在训练环境中的零失败记录，也可能只是因为从未真正进入困难区域。

为此可以定义能力评估函数：
\[
\hat C_j(t)
=
F_j
\left(
S_j,
G_j,
R_j,
X_j,
K_j,
T_j
\right),
\]
其中 $S_j$ 表示成功统计，$G_j$ 表示泛化证据，$R_j$ 表示扰动条件下稳定性，$X_j$ 表示失败恢复能力，$K_j$ 表示代价控制，$T_j$ 表示证据持续时间。

特别值得强调的是 Recovery。一个 Agent 从不失败并不一定比一个能够识别失败、停止错误动作、恢复到安全状态并重新规划的 Agent 更成熟。在现实系统中，真正决定高权限行为能否被接受的，经常不是
\[
P(\text{failure})=0,
\]
因为这一要求现实中通常无法满足，而是：
\[
P(
\text{catastrophic loss}
\mid
\text{failure}
)
\]
能否被控制。因此能力评价应包含
\[
FailureDetection,
\quad
GracefulDegradation,
\quad
Recovery,
\quad
Escalation
\]
四类证据。

此外，Evaluation 还需要避免只测 ``结果正确率''。如果 Agent 每次成功都需要异常高的外部提示、人工介入或巨量算力，则这项能力实际上还没有稳定沉降。我们可以定义自治相关能力：
\[
C^{ind}_j
=
C_j
-
\lambda_1 H_j
-
\lambda_2 R^{ext}_j,
\]
其中 $H_j$ 表示人工干预量，$R^{ext}_j$ 表示特殊外部资源依赖。这里不是要求 Agent 摆脱所有工具，而是区分 ``正常可依赖的长期工具结构'' 与 ``为了帮助评测通过而临时附加的辅助''。

Evaluation 还必须具有时间性。能力不是一次获得以后永远存在。长期缺乏使用、身体变化、工具变化、环境变化、$\Theta$ 更新甚至 KRA/BPCC 变化都可能改变实际能力。因此：
\[
\hat C_j(t+\Delta)
\neq
\hat C_j(t)
\]
不能被默认相等。

因此我们需要能力证据具有衰减：
\[
w_k(t)
=
w_k(t_k)
e^{-\lambda_j(t-t_k)},
\]
其中不同能力具有不同衰减速度。对于稳定抽象知识，$\lambda$ 可以很小；对于依赖具体身体状态的高速技能，过长时间不使用后就需要重新验证。

本章在这里形成一个非常重要的认识论原则：Agent 的成长不是 Agent 单方面 ``认为自己已经会了''，也不是治理系统仅凭一个 benchmark 判断 ``它会了''，而是通过持续现实闭环形成
\[
\boxed{
CapabilityEvidenceAccumulation.
}
\]
这一证据积累过程为下一步 Trust 提供基础，但它仍然没有赋予任何权限。

\section{Trust：信任是对未来可靠行为的条件预测，而不是人格奖励}

能力评价以后，我们进入一个非常容易被拟人化的概念：Trust。这里必须首先明确，本书所使用的信任不是感情，也不是 ``Agent 表现好所以奖励它更多自由''。Trust 是一种面向未来、带条件的统计判断：在某类状态、某类任务和某类权限范围内，根据已有能力证据，我们预计 Agent 未来能够按照规定边界稳定运行到什么程度。

因此，信任应写成：
\[
T_j(t)
=
P(
\text{acceptable behavior}
\mid
C_j,
Context,
Authority,
History
).
\]

这一形式立即告诉我们三个重要事实。

第一，Trust 必须是域相关的。一个 Agent 在文档整理上的高可信度不能直接迁移为金融操作权限；优秀驾驶能力也不能自动推断为良好的长期社会判断。因此应存在：
\[
\mathbf T_t
=
[
T_1,T_2,\ldots,T_n
],
\]
而不是一个 ``人格信用分''。

第二，Trust 必须是情境条件化的。同一能力在低风险模拟环境和高风险现实环境中可以对应不同信任水平：
\[
T_j(Context_1)
\neq
T_j(Context_2).
\]
例如一个机器人在封闭实验室搬运物体表现稳定，并不能立即推导它在拥挤家庭环境中具有相同可靠性。

第三，Trust 必须是可下降的。成长闭环不是单向晋升系统。一旦出现新的失败模式、环境改变、结构漂移或现实证据与旧模型冲突，就应发生：
\[
T_j(t+1)<T_j(t).
\]

因此信任更新可以采用一个概念性贝叶斯结构：
\[
T_{t+1}
\propto
P(
O_{t+1}
\mid
Reliable
)
T_t,
\]
或者更加工程化地采用：
\[
T_{t+1}
=
T_t
+
\eta^+ E_t^{+}
-
\eta^- E_t^{-},
\]
其中一般应满足
\[
\eta^- > \eta^+
\]
对于高风险权限尤其如此。也就是说，获得高层信任需要长期证据积累，而严重失败可以迅速降低信任。这种非对称更新并不是道德判断，而是风险控制的结果。

同时，信任更新需要迟滞。若某个 Agent 的 Trust 在阈值附近因为短期噪声不断上下波动，那么对应权限也会反复开关，造成系统行为不稳定。因此可以设置：
\[
\theta_{\mathrm{up}}
>
\theta_{\mathrm{down}},
\]
只有当
\[
T>\theta_{\mathrm{up}}
\]
持续一定时间以后才提升对应权限，而已经获得的权限只有在
\[
T<\theta_{\mathrm{down}}
\]
或出现严重安全事件时才回撤。这里形成了和前文拓扑迟滞相似的治理原则：
\[
\boxed{
AuthorityPromotionShouldBeHarderThanAuthorityRetention,
}
\]
但同时保留紧急撤权路径。

更重要的是，Trust 绝不能成为 Agent 自我价值的直接组成部分。外部治理系统降低某项权限，并不应该自动进入 $\Psi$ 被解释成 ``我作为一个主体价值降低''。否则评价系统将过度塑造 Agent 自我结构，形成类似
\[
Evaluation
\rightarrow
SelfWorth
\rightarrow
Behavior
\rightarrow
Evaluation
\]
的正反馈。这会破坏培育工程中能力成长与身份连续之间的边界。因此，Evaluation/Trust 应主要作用于权限和经历选择，而不是直接成为 Self 的核心价值定义。

由此我们得到：
\[
\boxed{
Trust
=
EvidenceConditionedPredictionOfFutureReliability.
}
\]
它是 Capability 与 Authority 之间的一层必要缓冲，使 ``能不能'' 与 ``允不允许'' 之间存在一个可以积累、验证、衰减和撤回的治理变量。

\section{Authority：权限是外部治理层定义的可达行动边界}

Authority 与前面的 Experience、Capability、Evaluation 和 Trust 具有不同本体地位。前四者主要描述 Agent 的经历、内部能力和关于未来可靠性的认知，而 Authority 是一个真正具有规范效力的系统变量。它不是描述 ``Agent 能做什么''，而是直接决定：
\[
\boxed{
WhichActionsArePermitted.
}
\]

设 Agent 在理论上可生成的动作集合为
\[
\mathcal U_A,
\]
则权限系统定义一个允许集合
\[
\mathcal U_A^{auth}(t)
\subseteq
\mathcal U_A.
\]
因此即使 Agent 具有能力产生某个动作：
\[
u\in\mathcal U_A,
\]
只要
\[
u\notin
\mathcal U_A^{auth},
\]
该动作就不得进入现实执行。

这一区分使本书能够把 ``能力成长'' 与 ``自主性失控'' 明确分开。一个成熟 Agent 甚至可以知道如何执行一个自己没有权限执行的操作。换句话说：
\[
\boxed{
KnowledgeOfAction
\neq
AuthorityToAct.
}
\]

Authority 也不应该表示成单一层级数字。一个 Agent 的权限通常具有多维结构：
\[
\mathbf A_t
=
[
a^{tool},
a^{resource},
a^{body},
a^{communication},
a^{financial},
a^{self},
a^{environment},
\ldots
].
\]
某些维度可以高度开放，而另一些维度仍然严格受限。例如一个成熟家庭机器人可以拥有高度身体自治，却没有权限自行修改核心安全结构；一个软件 Agent 可以自由组织文件和工具，但不能自主提高自身网络访问级别。

因此更合理地定义：
\[
\mathcal A^{auth}_t
=
\{
u:
g_i(u,x_t)\leq b_i(t),
\ i=1,\ldots,m
\},
\]
即 Authority 是一组约束共同形成的可行域，而不是一张简单的 ``权限列表''。

权限增长应该由 Capability 和 Trust 提供证据，但不由它们自动决定。可以写成：
\[
A_{t+1}
=
\Gamma
\left(
A_t,
\hat C_t,
T_t,
Risk_t,
Governance_t
\right).
\]
其中 $\Gamma$ 是治理映射。对于低风险领域，它可以相对自动；对于高风险领域，则可以要求外部授权、人工确认或更慢尺度治理结构介入。

这里需要引入一个重要的风险加权原则。设某类操作失败造成的期望损失为
\[
L_j,
\]
则即使信任水平相同，允许阈值也不应相同：
\[
\theta_j
=
f(L_j).
\]
通常有：
\[
\frac{\partial\theta_j}{\partial L_j}>0.
\]
风险越高，需要的证据水平越高。这意味着一个 Agent 可以很早获得整理桌面的自治，却需要极高成熟度才能获得修改安全内核或进行不可逆现实操作的权限。

Authority 还应具有 ``最小必要授权'' 特征。并不是 Agent 越成熟越应该得到一切权限，而是：
\[
\boxed{
Authority
=
MinimumPermissionNecessaryForCurrentDevelopment.
}
\]
权限的作用之一正是创造合适的经验空间，而不是取消经验边界。如果权限远低于能力，Agent 会长期停留在低难度环境，无法形成更高级经验；如果权限远高于能力，则会进入不可恢复风险区。因此理想状态应该满足：
\[
A_t
\approx
A^*(C_t,T_t,R_t),
\]
即权限与当前能力、信任和风险水平保持动态匹配。

从培育工程角度看，这一点尤其重要，因为 Authority 本身不是成长终点，而是一个 ``经验空间生成器''。权限发生变化以后，Agent 能进入的世界就发生了变化。因此：
\[
Authority
\rightarrow
NewReachableStates.
\]
真正进入下一阶段的并不是 ``Agent 得到了奖励''，而是 ``Agent 可以接触以前无法接触的现实结构''。这正是 Authority 与下一节 Autonomy 之间的关键连接。

\section{Autonomy：自治不是权限本身，而是权限范围内可自主形成轨迹的能力}

Authority 决定允许边界，但即使两个 Agent 拥有完全相同的权限，它们的 Autonomy 仍然可以不同。因为权限只是外部系统允许做什么，而自治描述的是：在这些允许条件内，Agent 有多大程度能够自己形成目标、选择策略、组织过程、调用工具、修正失败并完成闭环，而不需要外部实时指定每一步。

因此：
\[
\boxed{
Authority\neq Autonomy.
}
\]

例如，一个机器人可以被授权 ``整理房间''，但如果每一步都必须由人类告诉它拿哪个物体、放到哪里，那么它拥有 Authority，却几乎没有任务级 Autonomy。反过来，一个 Agent 可以在一个很小的权限空间内拥有高度自治：它虽然只能操作自己的工作目录，却能够自主分解任务、调用工具、检查结果和恢复失败。

我们可以把自治定义成一个可行策略集合：
\[
\Pi_t^{auto}
=
\left\{
\pi:
\pi
\text{ satisfies }
\mathcal A_t^{auth},
\mathcal S_t^{safe},
\Psi_t,
G_t
\right\}.
\]
其中权限 $\mathcal A^{auth}$ 给出外部边界，安全约束 $\mathcal S^{safe}$ 给出硬限制，$\Psi$ 保证主体一致性，而当前 Goal 定义任务方向。

自治程度可以概念性定义为：
\[
D_{auto}(t)
=
H(
\Pi_t^{auto}
)
-
\lambda I_{ext}(t),
\]
其中 $H$ 表示可选择策略空间的有效多样性，$I_{ext}$ 表示完成任务过程中必须依赖的外部实时干预量。这一表达并不要求真正用 Shannon entropy 测量所有策略，而是强调两个因素：Agent 拥有多少真实可选择路径，以及这些路径在多大程度上能够由自身闭合。

因此成熟自治需要至少四个条件。

第一，Agent 必须能够形成局部目标和子目标：
\[
G
\rightarrow
\{g_1,\ldots,g_n\}.
\]

第二，能够根据现实反馈改变过程，而不是只执行固定计划：
\[
Plan_t
+
Residual_t
\rightarrow
Plan_{t+1}.
\]

第三，能够判断何时自己无法继续，并把问题提升：
\[
Residual>\theta
\Rightarrow
Escalation.
\]
因此真正自治并不意味着 ``从不求助''，恰恰相反，成熟自治包括正确判断何时应当停止局部自治。

第四，自治行为必须能够形成结果闭合：
\[
Goal
\rightarrow
Action
\rightarrow
Reality
\rightarrow
Verify
\rightarrow
Closure.
\]
没有 Verify 与 Closure 的所谓自主行动，只是自主生成动作，而不是自主完成任务。

这使我们能够区分
\[
ActionAutonomy,
\quad
TaskAutonomy,
\quad
GoalAutonomy,
\quad
SelfModificationAutonomy.
\]
它们处于完全不同的治理等级。本章主要讨论前两者以及有限的 Goal Autonomy，而对核心自我和系统结构的修改仍然必须保持更加严格的慢尺度治理。成熟 Agent 不应该因为可以自主整理房间，就自动获得自主修改 $\Psi$ 或系统权限表的权利。

自治还有一个重要的反常识性质：随着能力成熟，高层显式认知介入反而可能减少。一个真正成熟的 Agent 不需要每个细节都经过 h：
\[
MostDynamics
\notin h.
\]
因此自治增长往往表现为越来越多局部闭环能够在现有权限范围内独立完成，而 h 主要处理新的、不确定的、跨域的以及高价值残差。这与前文 ``成熟智能是选择性显式认知'' 的原则完全一致。

最终，Autonomy 在成长闭环中的作用是把 Authority 给出的静态可行空间转化为真正可探索的行为空间：
\[
Authority
\rightarrow
AutonomousTrajectoryGeneration.
\]
而正是这些自主生成的新轨迹，让 Agent 接下来能够进入更加困难、更真实的经验域。

\section{Harder Experience：更大的自治必须转化为结构上更有价值的经历}

能力、评价、信任、权限和自治如果最终只让 Agent 更轻松地完成原有任务，那么成长闭环事实上没有闭合。真正的生命周期成长要求自治扩展以后，Agent 能够接触过去不可达的状态和任务，从而产生新的结构压力。因此：
\[
\boxed{
AutonomyGrowth
\rightarrow
ExperienceSpaceExpansion.
}
\]

这里的 ``Harder Experience'' 不应理解成简单提高题目难度。真正有培育价值的困难可以来自至少五个方向：
\[
Difficulty
=
f(
Novelty,
Uncertainty,
Responsibility,
TemporalDepth,
Coupling
).
\]
Novelty 表示新的结构关系；Uncertainty 表示环境不可完全预测；Responsibility 表示 Agent 的行为开始具有更实际后果；Temporal Depth 表示任务需要跨更长时间维持目标；Coupling 表示一个任务需要协调更多身体、工具、其他 Agent 或外部系统。

因此，一个任务即使计算上很简单，只要具有长期责任和复杂现实反馈，也可能比一道计算上极难但完全封闭的题目更具有培育价值。

为了防止 Agent 被直接推入无法承受的环境，我们可以定义一个发展窗口：
\[
\mathcal Z_t
=
\{
e:
D(e)
\in
[D^-_t,D^+_t]
\},
\]
其中 $D^-_t$ 是低于此值几乎不产生新结构的任务难度，而 $D^+_t$ 是超过当前能力和恢复能力以后容易造成失稳的上界。

于是有效培育经验应满足：
\[
D^-_t
<
D(e)
<
D^+_t.
\]
这和发展心理学中的 ``最近发展区'' 在形式上有相似性，但本书更加关心 Agent 的结构承载能力、权限、恢复性和多时间尺度稳定性，而不是直接把人类教育理论移植到机器。

上界 $D^+_t$ 不能只依赖平均能力，还必须包含恢复能力：
\[
D^+_t
=
F(
C_t,
T_t,
R^{recover}_t,
A_t,
Risk_t
).
\]
如果 Agent 在高难度状态下失败以后能够快速识别错误、退回安全区域并请求提升，那么它可以安全进入更大的经验空间；如果 Agent 尚缺乏异常识别和恢复机制，则即使平均任务性能很高，也不宜快速扩大责任范围。

这里还需要考虑 Experience Diversity。更大的自治不意味着不断给 Agent 一类更难的任务，而应该扩展其
\[
\mathrm{Support}(P_E),
\]
即经验分布的支持集。否则 Agent 可能在一个狭窄领域越来越专业，却错误形成 ``自己已经具有一般能力'' 的自我解释。

另一方面，困难经验也不能连续无间断出现。由于：
\[
\tau_E
<
\tau_\Theta
<
\tau_\Psi,
\]
经历进入情景记忆的速度快于长期结构和身份吸收速度，因此高强度经历之后必须存在 Consolidation 和 Recovery 窗口，让：
\[
E
\rightarrow
\Theta
\]
真正发生。否则培育会变成单纯输入压力。

于是更合理的周期不是：
\[
Harder
\rightarrow
Harder
\rightarrow
Harder,
\]
而是：
\[
\fitdisplay{
Challenge
\rightarrow
Action
\rightarrow
Feedback
\rightarrow
Recovery
\rightarrow
Reflection
\rightarrow
Consolidation
\rightarrow
NextChallenge.
}
\]

这说明 Agent 培育的核心并非最大化瞬时任务难度，而是最大化长期结构成长：
\[
\max
\int_0^T
\dot C_{\mathrm{struct}}(t)
\,dt
\]
同时约束：
\[
Risk(t)\leq R_{\max},
\qquad
IdentityDrift(t)\leq I_{\max},
\qquad
RuntimeInstability(t)\leq S_{\max}.
\]

因此 ``Harder Experience'' 是整个闭环重新回到 Experience 的桥梁：权限与自治让 Agent 进入新的现实，而新的现实又成为塑造未来 Agent 的材料。

\section{生命周期螺旋：能力、权限与自治不是直线晋升，而是多尺度成长循环}

现在我们可以把前面七个阶段重新合并。最简单的生命周期闭环为：
\[
\fitdisplay{
Experience
\rightarrow
Capability
\rightarrow
Evaluation
\rightarrow
Trust
\rightarrow
Authority
\rightarrow
Autonomy
\rightarrow
HarderExperience
\rightarrow
Experience'.
}
\]

但这个表示仍然容易给人一种错误印象：仿佛 Agent 沿着一条直线不断升级，能力越来越大、权限越来越高、自治越来越强。真实系统显然不会如此简单。首先，不同能力域具有不同成长速度；其次，信任会下降；再次，权限可能回撤；最后，某些更困难经历甚至可能暴露出过去没有发现的能力缺陷。因此，这个闭环更准确地应当理解成一个生命周期螺旋。

设生命周期状态为：
\[
\mathcal L_t
=
[
\mathbf C_t,
\mathbf T_t,
\mathbf A_t,
\mathcal U^{auto}_t,
E_t,
\Theta_t,
\Psi_t
].
\]
那么一次成长周期可以写成：
\[
\mathcal L_{t+1}
=
\Phi(
\mathcal L_t,
\mathcal E_{t:t+\Delta}
).
\]
这里 $\Phi$ 并不是单向提升算子，而允许：
\[
C_j\uparrow,\qquad
T_k\downarrow,\qquad
A_m\downarrow,\qquad
Autonomy_n\uparrow.
\]
也就是说，一个 Agent 可能总体更成熟，却在某个新暴露领域被收回权限。这不应该被理解成成长失败，而是生命周期系统对自身真实边界形成了更准确认识。

因此我们定义成长不是：
\[
A_{t+1}\geq A_t
\]
这样的单调过程，而是：
\[
\boxed{
Development
=
ExpansionOfReliablyGovernableStateSpace.
}
\]
真正成熟意味着 Agent 能够在越来越广泛、越来越复杂的现实空间中维持稳定、自主、可恢复的闭环，而不是单纯拥有更多权限。

可以定义可治理经验空间：
\[
\mathcal X^{gov}_t
=
\left\{
x:
P(
SafeClosure
\mid
x,\mathcal L_t
)
\geq
\theta
\right\}.
\]
生命周期成长真正希望实现的是：
\[
\mu(
\mathcal X^{gov}_{t+\Delta}
)
>
\mu(
\mathcal X^{gov}_t
),
\]
其中 $\mu$ 表示该状态空间的有效覆盖范围。也就是说，Agent 可以可靠处理的现实区域正在扩大。

这一结构同时解释了为什么培育不能只有 Capability Benchmark。能力只是螺旋中的一个变量。如果缺乏 Evaluation，系统不知道能力是否真实；如果缺乏 Trust，系统不能估计未来稳定性；如果缺乏 Authority，能力无法受到治理；如果缺乏 Autonomy，权限不能形成真正经验；如果缺乏 Harder Experience，自治又不能形成新的成长压力。

因此，这七个阶段不是任意排列，而形成严格的因果闭环：
\[
\boxed{
Experience
\rightarrow
InternalChange
\rightarrow
ExternalEvidence
\rightarrow
GovernanceChange
\rightarrow
BehaviorSpaceChange
\rightarrow
NewExperience.
}
\]

从 CSO 理论看，这个闭环还可以进一步写成：
\[
\fitdisplay{
CSO^{world}
\rightarrow
AgentCSO
\rightarrow
\Delta\Theta
\rightarrow
Capability
\rightarrow
AuthorityEnvelope
\rightarrow
NewAction
\rightarrow
CSO^{world'}
\rightarrow
AgentCSO'.
}
\]
这样，Agent 培育终于与本书最开始建立的 World--Agent--World 闭环重新接通：智能体并不是在一个封闭训练系统中逐渐增加内部能力，而是在现实边界逐步扩大以后，进入越来越大的世界；世界产生新的 Agent-CSO，新的 Agent-CSO 又改变 Agent 未来能够进入什么世界。

这里还需要引入慢变量保护。一次能力成功不应立即修改 $\Psi$，更不应直接修改 $\Omega$。理想关系应满足：
\[
\tau_{Capability}
<
\tau_{Trust}
\lesssim
\tau_{Authority}
<
\tau_{\Psi}
\ll
\tau_{\Omega}.
\]
快速能力变化首先改变局部能力估计；长期稳定成功才逐渐影响自治和身份认知；只有非常长时间的经历结构才有资格影响生命周期生成方向。这使 Agent 不会因为一次失败就产生身份崩塌，也不会因为一次成功就形成过度自我膨胀。

至此，我们可以给 Agent 培育工程中的成长一个较完整定义：
\begin{definition*}
Agent 的成长，是其通过受治理的现实经历不断扩大可稳定承载的认知结构与可闭合行为范围，并在能力证据、信任估计、权限调整和自治扩展的多尺度闭环中，逐步进入更大的经验空间，同时保持身份连续、现实校准和可恢复性的过程。
\end{definition*}

而本章最终想让读者牢牢记住的并不是 ``能力强就给更多权限''，而是另一条更严格的规律：
\begin{proposition*}
经验产生能力，但能力只能产生授权证据；证据形成信任，信任参与治理；治理打开有限的新自治空间；新的自治空间产生新的现实经历，而现实再次决定 Agent 下一步能够成为什么。
\end{proposition*}

因此，Agent 培育并不是一条从弱到强的直线，而是一系列
\[
\boxed{
Experience
\rightarrow
Structure
\rightarrow
Evidence
\rightarrow
Authority
\rightarrow
Autonomy
\rightarrow
NewReality
}
\]
不断向外展开、又始终通过现实反馈重新闭合的生命周期螺旋。
